\begin{proof}

\begin{enumerate}[label={$\langle1\rangle$\arabic*.},leftmargin=*]
\item Put $h=x^a-y^bz^c$. We first compare the actual map
with its reduced equation. This polynomial is weighted
homogeneous of degree $a(b+c)$ for weights
$(b+c,a,a)$. On a small simply connected coordinate
neighbourhood choose a holomorphic logarithm of $u$.
Define the holomorphic coordinate change
\[
 \Phi(x,y,z)=\left(
 x u^{1/(ma)},\;
 y u^{1/(m(b+c))},\;
 z u^{1/(m(b+c))}\right).
                                                        \tag{A.J.34}\hypertarget{tag-0065}{}\label{cdp:eq:A.J.34}
\]
Its derivative at the origin is an invertible diagonal
matrix. The holomorphic inverse function theorem and
weighted homogeneity give
$h\circ\Phi=u^{1/m}h$ and therefore
$f=(h\circ\Phi)^m$ exactly. For a sufficiently small
nonzero regular value $t$, a chosen Milnor fibre of
$h$ at one of the $m$ roots of $t$ is consequently a
compact oriented submanifold with boundary of the
smooth global fibre $Y_t$. Denote this Milnor fibre
by $M$.

\item We first give a direct spherical cycle, without an
irreducibility assumption. Put $d=\gcd(b,c)$,
$b'=b/d$ and $c'=c/d$, and consider the smooth affine
curve
\[
 B=\{(x,v)\in\C^2:x^a-v^d=1\}.
\]
The affine surface $H=\{h=1\}$ maps holomorphically
to $B$ by $v=y^{b'}z^{c'}$. The curve $B$ is
connected: projection to the $v$-line has $a$ sheets,
and a small loop about any root of $v^d+1$ cyclically
permutes them. It is smooth because its two partial
derivatives cannot vanish on $x^a-v^d=1$.
The function $v$ has $a$ distinct simple zeros
$P_j=(\zeta_a^j,0)$ on $B$.

\item Choose two of these zeros, $P$ and $Q$, and an
embedded arc $\gamma:[0,1]\to B$ joining them and
avoiding the other zeros. Choose smooth complex
functions $A(t),B_0(t)$ on its interior such that
\[
 A(t)^{b'}B_0(t)^{c'}=v(\gamma(t)),\qquad
 |A(t)|=|v(\gamma(t))|^{1/(2b')},\quad
 |B_0(t)|=|v(\gamma(t))|^{1/(2c')}.
\]
They exist by choosing a continuous argument along
the arc. The continuous map
\[
 (t,\theta)\longmapsto
 \bigl(x(\gamma(t)),A(t)e^{\mathrm i c'\theta},
                 B_0(t)e^{-\mathrm i b'\theta}\bigr)
                                                        \tag{A.J.33a}\hypertarget{tag-0063}{}\label{cdp:eq:A.J.33a}
\]
from $[0,1]\times S^1$ to $H$ collapses each boundary
circle to the corresponding point $P$ or $Q$.
It therefore defines an oriented spherical cycle
$\sigma_\gamma:S^2\to H$. Coprimality of $b'$ and
$c'$ makes each interior circle embedded. The whole
map is injective away from the collapsed ends,
because the base arc is embedded.

\item Let $\gamma'$ be a nearby arc with the same endpoints
and disjoint interior, obtained by an isotopy in a
thin disc neighbourhood of $\gamma$. The isotopy
keeps all other zeros of $v$ outside that disc.
Choose the arguments in \eqref{cdp:eq:A.J.33a} continuously during
the isotopy; both coordinates tend uniformly to zero
at the ends. This gives a continuous homotopy of
maps from $S^2$, so the resulting spherical cycles
are homologous. Their supports meet only at $P$ and
$Q$. Near either endpoint, $(y,z)$ are actual
holomorphic coordinates on $H$, since $a x^{a-1}$
is nonzero there. Arrange the two arcs to be radial
with different arguments in the local $v$-coordinate.
On a sufficiently small sphere in these coordinates,
the links of the two cycles are disjoint parallel
curves of slope $(c',-b')$ on a Clifford torus.
Their oriented linking number is $-b'c'$.
Indeed the torus framing of a $(p,q)$ curve has
linking number $pq$ with its parallel copy: push the
two curves to the opposite sides of the Heegaard
splitting of $S^3$ into its two solid tori and count
meridian discs, obtaining $p q$ signed intersections.
Here $p=c'$ and $q=-b'$.

\item The local intersection of the two coned cycles is
that linking number. This equality follows directly
by replacing the two cones inside the small ball by
transverse relative two-chains with the same boundary
links; their intersection is the defining linking
number in its boundary three-sphere. At $Q$ both
cycle orientations reverse relative to the outgoing
radial coordinate, so the two reversals leave the
local intersection sign unchanged. Consequently
\[
 [\sigma_\gamma]^2
   =[\sigma_\gamma]\cdot[\sigma_{\gamma'}]
   =-2b'c'\ne0.                                  \tag{A.J.33b}\hypertarget{tag-0064}{}\label{cdp:eq:A.J.33b}
\]
Both cycles and their homotopy have compact support.
The weighted homothety
$(x,y,z)\mapsto(s^{b+c}x,s^a y,s^a z)$ maps
$H$ biholomorphically onto $\{h=s^{a(b+c)}\}$.
For sufficiently small positive $s$, it puts those
compact supports into any prescribed Milnor
neighbourhood of the origin. Thus \eqref{cdp:eq:A.J.33b} is a
spherical class of nonzero square in an actual
local Milnor fibre $M$.

\item For an independent check when $\gcd(a,b,c)=1$,
we compute the full homotopy type and the rank of
the intersection form from two primary results.
The join theorem
\cite[Theorem L1 and Example 2.4(a)]{Nemethi1991Join}
applies to arbitrary analytic germs, including
non-isolated ones. Set $d=\gcd(b,c)$. The Milnor
fibre of $y^bz^c$ is a disjoint union of $d$ annuli:
its equations in a Milnor polydisc are
$y^{b/d}z^{c/d}=\zeta\epsilon^{1/d}$, and each
component is parametrized by
$y=A w^{c/d}$, $z=B w^{-b/d}$. Coprimality makes
this parametrization injective, and the two radius
bounds give a nonempty annulus in $w$.
The Milnor fibre of $x^a$ consists of $a$ points.
The join theorem thus identifies $M$, by a homotopy
equivalence, with the union of $a$ cones on a
disjoint union of $d$ circles, all glued along their
common base. Contracting one cone and then choosing
arcs between the $d$ base components gives
\[
 M\simeq
 \bigvee^{(a-1)d}S^2\ \vee\
 \bigvee^{(a-1)(d-1)}S^1.                          \tag{A.J.35}\hypertarget{tag-0066}{}\label{cdp:eq:A.J.35}
\]
In particular, every class in $H_2(M,\Z)$ is a sum
of images of spherical fundamental classes.

\item For the same irreducible germ, the boundary calculation
of \cite[Theorem 4.1 and Lemma 4.6]{MPW2005Hirzebruch}
gives a Seifert three-manifold with base genus
\[
 g_\partial=(a-1)(d-1),\qquad
 e_0=\frac{a}{(b/d)(c/d)}>0.                     \tag{A.J.36}\hypertarget{tag-0067}{}\label{cdp:eq:A.J.36}
\]
Here $e_0$ is its rational Seifert Euler number in
the orientation convention of that paper. Its
nonzero value gives
$b_1(\partial M)=2g_\partial$. To make this last
deduction explicit, the rational abelianized Seifert
presentation has a regular-fibre generator $v$,
exceptional-fibre generators $q_i$, the relations
$\alpha_iq_i+\beta_iv=0$ and
$\sum_iq_i+ev=0$, and $2g_\partial$ free base
generators. Eliminating the $q_i$ gives $e_0v=0$.
Thus all fibre generators vanish over $\Q$, leaving
exactly the $2g_\partial$ base generators.

\item Let $I_M:H_2(M,\Q)\to H_2(M,\partial M;\Q)$ be
the intersection morphism. Poincare--Lefschetz duality
and the homology sequence of the pair give
\[
 \operatorname{rank} I_M
   =b_2(M)+b_1(M)-b_1(\partial M)=a-1>0.          \tag{A.J.37}\hypertarget{tag-0068}{}\label{cdp:eq:A.J.37}
\]
Indeed $H_3(M)=0$ by \eqref{cdp:eq:A.J.35}, so
$H_3(M,\partial M)\simeq H^1(M)$ injects into
$H_2(\partial M)$. Its cokernel is the kernel of
$I_M$, which has dimension
$b_1(\partial M)-b_1(M)$. This proves the displayed
rank formula. For $a=2$, $b=1$, $c=\ell$, the
more precise boundary is $L(2\ell,1)$ by
\cite[Theorem 4.2(2)]{MPW2005Hirzebruch}; the same
pair sequence over $\Z$ makes the rank-one
intersection form have absolute determinant
$2\ell$. The special case $\ell=1$ is the ordinary
quadratic singularity and has determinant two.

\item Finally let $j:M\hookrightarrow Y_t$ be the actual
embedding obtained from \eqref{cdp:eq:A.J.34}. A complex two-torus
has contractible universal cover $\C^2$, so
$\pi_2(Y_t)=0$. In particular the spherical class
constructed in \eqref{cdp:eq:A.J.33a} maps to zero in
$H_2(Y_t,\Z)$. Intersection numbers of compact
two-cycles in the interior of $M$ are preserved by
this oriented codimension-zero embedding: the same
small transverse perturbations compute them in
both manifolds. Its square in $Y_t$ is therefore
both zero and $-2b'c'$, contradicting \eqref{cdp:eq:A.J.33b}.
In the irreducible case, \eqref{cdp:eq:A.J.35}--\eqref{cdp:eq:A.J.37} provide
an additional contradiction on the whole intersection
space.
This proves the exclusion for the actual map,
including every multiplicity in \eqref{cdp:eq:A.J.33}.

\end{enumerate}
\end{proof}
